\subsection{问题定义与总体目标}
\label{sec:purpose}

本课题面向低时延 LLM 推理，旨在构建\textbf{依赖感知的 Megakernel 推理融合框架与系统}。输入是任意计算图；框架首先通过依赖分析与任务拆解，为图中每条相邻算子边选择合适的执行方式，并生成可落地的端到端执行计划。核心问题因此包含前后衔接的两层决策：其一，依据 task / tile 级依赖形态，明确各边融合方案并在全图上搜索最优边处理方案；其二，对跨 grid 边进一步以 PDL 组织依赖感知的跨边界重叠。

形式化地，设输入计算图为 $G=(O,E)$，其中 $O$ 为算子集合，$E$ 为算子间数据依赖。框架将每个算子 $o\in O$ 拆解为可调度的 task 集合 $\mathcal{T}(o)$；对相邻算子边 $e=(o_p,o_c)$，进一步得到从生产者 task 到消费者 task 的依赖关系。该关系刻画消费者需要等待哪些生产者 task、依赖按 task/tile 粒度局部就绪还是全量就绪，以及消费者在读取依赖数据前是否存在可独立执行的工作。

对每条边 $e$，框架生成调度候选集合 $\mathcal{M}(e)$，并将边处理方案表示为 $\pi(e)=\langle m_e,\theta_e\rangle$，其中 $m_e\in\mathcal{M}(e)$ 为处理范式，$\theta_e$ 为相应的同步、投放或 PDL 参数。本课题重点研究表~\ref{tab:edge-scheduling-modes} 中四类候选。前两类将边保留在同一 grid 内，分别利用局部就绪或全量就绪推进消费者；后两类将边映射为物理 grid 边界，并进一步决定是否使用 PDL。全图执行计划记为 $\Pi=\{\pi(e)\mid e\in E\}$。边的选择并非彼此独立：一条边是否跨 grid 会改变相邻算子的 grid 分组、资源状态和其它边的可重叠窗口，因此需要在全图上联合求解。

% \begin{table}[htbp]
%   \centering
%   \caption{计算图相邻算子边的候选处理范式}
%   \label{tab:edge-scheduling-modes}
%   \small
%   \begin{tabularx}{\textwidth}{p{0.22\textwidth}X}
%     \toprule
%     处理范式 & 适用依赖与执行方式 \\
%     \midrule
%     同 grid 细粒度事件
%       & 消费者 task 只依赖部分生产者 task；对应 task/tile 就绪后即可通过 event 或 counter 放行，并与生产者尾部重叠执行。 \\
%     同 grid 全量就绪
%       & 消费者主体依赖生产者的完整输出；消费者在内部 ready counter 或调度门控满足后开始，避免错误读取未完成数据。 \\
%     跨 grid completion
%       & 在该边形成物理边界；由生产者 grid completion 统一保证完成顺序与内存可见性，再调度消费者 grid。 \\
%     跨 grid PDL
%       & 保留物理边界，但在 programmatic completion 后提前调度消费者；消费者先执行不读取生产者输出的独立前缀，并在首次依赖读取前同步。 \\
%     \bottomrule
%   \end{tabularx}
% \end{table}

% \usetikzlibrary{arrows.meta,calc,positioning}
\begin{figure}[htbp]
  \centering
  \makebox[\textwidth][c]{%
  \resizebox{0.98\textwidth}{!}{%
  \begin{tikzpicture}[
      x=1cm,
      y=1cm,
      font=\sffamily\small,
      panel/.style={
        draw=black!24,
        fill=black!1,
        rounded corners=2.5pt,
        line width=0.55pt
      },
      paneltitle/.style={
        anchor=west,
        font=\sffamily\small\bfseries,
        text=black!78
      },
      op/.style={
        circle,
        draw=blue!62!black,
        fill=blue!10,
        minimum size=7mm,
        inner sep=0pt,
        line width=0.6pt
      },
      ptask/.style={
        draw=blue!62!black,
        fill=blue!10,
        rounded corners=1.2pt,
        minimum width=10mm,
        minimum height=5.5mm,
        inner sep=1pt
      },
      ctask/.style={
        draw=teal!62!black,
        fill=teal!10,
        rounded corners=1.2pt,
        minimum width=10mm,
        minimum height=5.5mm,
        inner sep=1pt
      },
      mode/.style={
        draw=violet!55!black,
        fill=violet!7,
        rounded corners=1.5pt,
        minimum width=19mm,
        minimum height=7mm,
        align=center,
        font=\sffamily\scriptsize
      },
      search/.style={
        draw=teal!58!black,
        fill=teal!7,
        rounded corners=1.8pt,
        minimum width=41mm,
        minimum height=8mm,
        align=center,
        font=\sffamily\scriptsize\bfseries
      },
      grid/.style={
        draw=black!34,
        fill=white,
        rounded corners=1.8pt,
        line width=0.55pt
      },
      tag/.style={
        draw=black!25,
        fill=white,
        rounded corners=1pt,
        inner xsep=3pt,
        inner ysep=1.5pt,
        font=\sffamily\scriptsize,
        text=black!68
      },
      flow/.style={
        -{Latex[length=1.8mm]},
        draw=black!55,
        line width=0.75pt
      },
      dep/.style={
        -{Latex[length=1.3mm]},
        draw=black!50,
        line width=0.5pt
      },
      chosen/.style={
        -{Latex[length=1.5mm]},
        draw=red!62!black,
        line width=0.8pt
      }
    ]

    % ------------------------------------------------------------------
    % Top row: input -> dependency analysis -> candidate/search
    % ------------------------------------------------------------------
    \draw[panel] (0,3.25) rectangle (4.05,7.35);
    \node[paneltitle] at (0.28,7.02) {1.\ 输入计算图};

    \node[op] (a) at (0.72,5.45) {$A$};
    \node[op] (b) at (1.78,6.10) {$B$};
    \node[op] (c) at (1.78,4.80) {$C$};
    \node[op] (d) at (3.02,5.45) {$D$};
    \node[op] (e) at (3.02,4.25) {$E$};
    \draw[dep] (a) -- (b);
    \draw[dep] (a) -- (c);
    \draw[chosen] (b) -- node[above, font=\sffamily\scriptsize,
      text=red!62!black] {$e$} (d);
    \draw[dep] (c) -- (d);
    \draw[dep] (c) -- (e);
    \node[font=\sffamily\scriptsize, text=black!58, align=center]
      at (2.02,3.62) {选择图中每条相邻算子边\\的执行方式};

    \draw[panel] (4.55,3.25) rectangle (10.25,7.35);
    \node[paneltitle] at (4.83,7.02) {2.\ task/tile 级依赖分析};

    \node[font=\sffamily\scriptsize\bfseries, text=blue!58!black]
      at (5.78,6.45) {生产者 tasks};
    \node[font=\sffamily\scriptsize\bfseries, text=teal!58!black]
      at (8.95,6.45) {消费者 tasks};
    \node[ptask] (p1) at (5.30,5.72) {$p_1$};
    \node[ptask] (p2) at (6.22,5.72) {$p_2$};
    \node[ptask] (p3) at (5.30,4.88) {$p_3$};
    \node[ptask] (p4) at (6.22,4.88) {$p_4$};
    \node[ctask] (q1) at (8.55,5.72) {$q_1$};
    \node[ctask] (q2) at (9.46,4.88) {$q_2$};
    \draw[dep] (p1) -- (q1);
    \draw[dep] (p2) -- (q1);
    \draw[dep] (p2) -- (q2);
    \draw[dep] (p3) -- (q2);
    \draw[dep] (p4) -- (q2);

    \node[tag] at (5.42,3.82) {就绪粒度};
    \node[tag] at (6.82,3.82) {读写集合};
    \node[tag] at (8.20,3.82) {资源占用};
    \node[tag] at (9.50,3.82) {独立前缀};

    \draw[panel] (10.75,3.25) rectangle (16.10,7.35);
    \node[paneltitle] at (11.03,7.02) {3.\ 逐边候选与全图搜索};

    \node[mode] (m1) at (12.05,6.12) {同 grid\\细粒度事件};
    \node[mode] (m2) at (14.78,6.12) {同 grid\\全量就绪};
    \node[mode] (m3) at (12.05,5.08) {跨 grid\\completion};
    \node[mode] (m4) at (14.78,5.08) {跨 grid\\PDL};
    \node[search] (s) at (13.415,3.90)
      {合法性剪枝 \quad 成本估计 \quad 全图联合搜索};
    \draw[flow] (m1.south) -- ++(0,-0.26) -| (s.north);
    \draw[flow] (m2.south) -- ++(0,-0.26) -| (s.north);
    \draw[flow] (m3.south) -- ++(0,-0.17) -| (s.north);
    \draw[flow] (m4.south) -- ++(0,-0.17) -| (s.north);

    \draw[flow] (4.05,5.30) -- (4.55,5.30);
    \draw[flow] (10.25,5.30) -- (10.75,5.30);

    % ------------------------------------------------------------------
    % Bottom row: generated physical plan
    % ------------------------------------------------------------------
    \draw[panel] (0,0) rectangle (16.10,2.68);
    \node[paneltitle] at (0.28,2.37)
      {4.\ 输出：覆盖全图的物理执行计划};

    \draw[grid] (0.50,0.45) rectangle (7.10,1.93);
    \node[font=\sffamily\scriptsize\bfseries, text=black!65]
      at (1.22,1.68) {Grid 0};
    \node[ptask, minimum width=13mm] (ga) at (2.30,1.12) {task $A$};
    \node[ctask, minimum width=13mm] (gb) at (4.30,1.12) {task $B$};
    \node[ctask, minimum width=13mm] (gc) at (6.25,1.12) {task $C$};
    \draw[chosen] (ga) -- node[above, font=\sffamily\scriptsize,
      text=red!62!black] {event} (gb);
    \draw[dep] (gb) -- node[above, font=\sffamily\scriptsize,
      text=black!58] {counter} (gc);

    \fill[orange!8] (7.42,0.40) rectangle (7.74,1.98);
    \draw[orange!72!black, densely dashed, line width=0.75pt]
      (7.58,0.40) -- (7.58,1.98);
    \node[font=\sffamily\tiny, text=orange!72!black]
      at (7.58,0.20) {grid boundary};

    \draw[grid] (8.05,0.45) rectangle (15.62,1.93);
    \node[font=\sffamily\scriptsize\bfseries, text=black!65]
      at (8.78,1.68) {Grid 1};
    \node[ctask, minimum width=15mm] (prefix) at (9.60,1.12) {prefix};
    \node[
      draw=orange!72!black,
      fill=orange!10,
      densely dotted,
      rounded corners=1.2pt,
      minimum width=12mm,
      minimum height=5.5mm,
      inner sep=1pt
    ] (wait) at (12.02,1.12) {wait};
    \node[ctask, minimum width=17mm] (body) at (14.35,1.12) {dependent body};
    \draw[dep] (prefix) -- (wait);
    \draw[dep] (wait) -- (body);
    \draw[chosen, dashed] (6.95,1.68) to[out=15,in=165]
      node[above, font=\sffamily\scriptsize, text=red!62!black]
      {PDL trigger} (8.18,1.68);

    \draw[flow] (13.42,3.25) -- (13.42,2.68);

  \end{tikzpicture}%
  }}%
  \caption{依赖感知的计算图边调度框架。框架从任意计算图出发，在 task/tile 粒度分析每条边的就绪关系、读写集合、资源占用与独立前缀；随后为各边生成候选并进行全图联合搜索，最终输出 grid 分组、边同步方式和 PDL 配置。}
  \label{fig:edge-scheduling-framework}
\end{figure}


基于以上统一依赖模型，总体优化目标为
\begin{equation}
  \Pi^\star(G,x)
  =\arg\min_{\Pi\in\mathcal{F}(G,x)}
    T_{\mathrm{decode}}(G,x;\Pi),
  \label{eq:global-scheduling-objective}
\end{equation}
其中 $x$ 表示输入规模与运行配置，$\mathcal{F}(G,x)$ 表示满足数据依赖、内存可见性和硬件资源约束的合法计划集合，$T_{\mathrm{decode}}$ 表示完整解码路径的端到端时延。式~\eqref{eq:global-scheduling-objective} 将每条边的处理范式及参数作为全图决策变量，联合确定 grid 分组、同步方式与 PDL 配置。




\subsection{主要研究内容}

围绕上述目标，本课题重点研究以下四个问题。

\begin{enumerate}
  \item \textbf{计算图依赖关系建模。}
  建立从算子图到 task graph 的 lowering 过程，以统一抽象刻画相邻算子边的 task/tile 级依赖关系，并据此构造各边合法的处理候选空间。

  \item \textbf{依赖感知的 PDL 跨边界重叠。}
  在依赖抽象之上，研究跨 grid 边的 PDL 变换：依据就绪关系识别可提前发起的重叠工作，在保持同步合法的前提下，确定重叠范围、发起时机与执行形态，使生产者与消费者之间的可重叠部分尽可能隐藏边界等待并优化数据读取。

  \item \textbf{全图边处理方案搜索。}
  面向全图构造边处理与 grid 分组的联合搜索问题；以静态分析、轻量运行时特征与性能模型估计各候选的收益与开销，在合法空间内求解式~\eqref{eq:global-scheduling-objective}，得到端到端时延最优的执行计划。

  \item \textbf{框架实现与端到端验证。}
  实现计划生成、CUDA 代码生成与运行时执行，将所得方案接入 decoder layer、完整 decode step 和连续 token 生成；真实的推理workload中进行精度验证的和性能测试。
\end{enumerate}

\subsection{预期目标、创新点与成果}

\subsubsection{预期目标}

与主要研究内容一一对应，本课题预期达到：

\begin{enumerate}
  \item 完成计算图依赖关系建模，形成从算子图到 task graph 的统一抽象，并为相邻算子边给出合法的处理候选空间。
  \item 完成依赖感知的 PDL 跨边界重叠方法，能够依据就绪关系识别重叠工作，并在同步合法约束下确定跨 grid 边的重叠方案。
  \item 完成全图边处理与 grid 分组的联合搜索，在合法空间内求解式~\eqref{eq:global-scheduling-objective}，得到端到端时延更优的执行计划。
  \item 完成框架原型与端到端验证，贯通计划生成、CUDA 代码生成与运行时执行，并在真实推理 workload 上完成精度与性能评估。
\end{enumerate}

\subsubsection{预期创新点}

\begin{enumerate}
  \item 提出面向任意计算图的依赖感知边处理抽象，以 task/tile 级依赖关系统一刻画同 grid 与跨 grid 候选，形成可搜索的合法方案空间。
  \item 提出依赖感知的 PDL 跨边界重叠方法，将重叠工作识别、发起时机与执行形态纳入统一变换，而不仅是固定边界上的静态预取。
  \item 提出覆盖计算图全部边的联合执行计划搜索方法，使边处理选择与 grid 分组协同优化，直接服务于端到端 decode 时延目标。
\end{enumerate}

\subsubsection{预期成果}

预期形成依赖感知的 Megakernel 推理融合框架与系统原型，覆盖依赖建模、PDL 跨边界重叠、全图计划搜索及代码生成运行时；并给出真实推理 workload 上的精度与性能实验结果，以及硕士学位论文和相关学术论文。
